\documentclass[12pt]{article}

% Packages
\usepackage[a4paper, margin=1in]{geometry}
\usepackage{setspace}
\usepackage{amsmath, amssymb}
\usepackage{hyperref}
\usepackage{titlesec}
\usepackage{enumitem}

% Formatting
\onehalfspacing
\setlist[itemize]{noitemsep, topsep=0pt}
\setlist[enumerate]{noitemsep, topsep=0pt}

\titleformat{\section}{\large\bfseries}{\thesection.}{0.5em}{}
\titleformat{\subsection}{\normalsize\bfseries}{\thesubsection}{0.5em}{}
\titleformat{\subsubsection}{\normalsize\itshape}{\thesubsubsection}{0.5em}{}

\begin{document}

\begin{center}
    {\Large \textbf{CSE400 -- Fundamentals of Probability in Computing}}\\[0.2cm]
    {\large Lecture 7: Expectation, CDFs, PDFs and Problem Solving}\\[0.3cm]
    \textbf{Instructor:} Dhaval Patel, PhD\\
    \textbf{Date:} January 27, 2025
\end{center}


\section{The Cumulative Distribution Function (CDF)}

\subsection{Definition}

Let $X$ be a random variable. The \textbf{Cumulative Distribution Function (CDF)} of $X$ is defined as
\[
F_X(x) = P(X \le x).
\]

\subsubsection*{Step-by-Step Reasoning}

\begin{enumerate}
    \item Consider the event $\{X \le x\}$.
    \item This event contains all outcomes for which the random variable takes a value less than or equal to $x$.
    \item The probability of this event defines the cumulative distribution up to the point $x$.
\end{enumerate}

This definition applies to both discrete and continuous random variables.

\subsection{Properties of the CDF}

From the definition $F_X(x) = P(X \le x)$, the following properties follow directly:

\begin{enumerate}
    \item \textbf{Non-decreasing property}\\
    If $x_1 < x_2$, then $\{X \le x_1\} \subseteq \{X \le x_2\}$, and hence
    \[
    F_X(x_1) \le F_X(x_2).
    \]

    \item \textbf{Boundary conditions}
    \[
    \lim_{x \to -\infty} F_X(x) = 0, \qquad
    \lim_{x \to +\infty} F_X(x) = 1.
    \]

    \item \textbf{Right-continuity}\\
    The CDF is right-continuous due to properties of probability measures.

    \item \textbf{Interval probability}
    For $a < b$,
    \[
    P(a < X \le b) = F_X(b) - F_X(a).
    \]
\end{enumerate}

\subsection{Example (as per slide structure)}

An example CDF is obtained by accumulating probabilities for a discrete random variable or by integrating a density function for a continuous random variable up to the value $x$. No additional numerical example is introduced.

\section{The Probability Density Function (PDF)}

\subsection{Definition}

For a \textbf{continuous random variable} $X$, the \textbf{Probability Density Function (PDF)} $f_X(x)$ is defined through its relationship with the CDF:
\[
F_X(x) = \int_{-\infty}^{x} f_X(t)\, dt.
\]

Where differentiable,
\[
f_X(x) = \frac{d}{dx} F_X(x).
\]

\subsubsection*{Step-by-Step Reasoning}

\begin{enumerate}
    \item The CDF accumulates probability up to $x$.
    \item The PDF represents the rate of change of this accumulated probability.
    \item Differentiation of the CDF yields the PDF.
\end{enumerate}

\subsection{Properties of the PDF}

\begin{enumerate}
    \item \textbf{Non-negativity}
    \[
    f_X(x) \ge 0 \quad \text{for all } x.
    \]

    \item \textbf{Total probability equals one}
    \[
    \int_{-\infty}^{+\infty} f_X(x)\, dx = 1.
    \]

    \item \textbf{Probability over an interval}\\
    For $a < b$,
    \[
    P(a < X < b) = \int_a^b f_X(x)\, dx.
    \]
\end{enumerate}

\subsection{Example (as per slide structure)}

Given a PDF, probabilities are computed by integrating over the desired interval, and the CDF is obtained by integrating from $-\infty$ to $x$.

\section{Expectation of Random Variables}

\subsection{Definition}

The \textbf{Expectation} of a random variable $X$ is defined as:

\begin{itemize}
    \item \textbf{Discrete case}
    \[
    \mathbb{E}[X] = \sum_x x \, P(X = x).
    \]

    \item \textbf{Continuous case}
    \[
    \mathbb{E}[X] = \int_{-\infty}^{+\infty} x \, f_X(x)\, dx.
    \]
\end{itemize}

\subsubsection*{Step-by-Step Reasoning}

\begin{enumerate}
    \item Each value of $X$ is weighted by its probability or density.
    \item The sum or integral aggregates these weighted values.
    \item The result represents the long-run average value.
\end{enumerate}

\subsection{Example (as per slide structure)}

Expectation is computed directly using the PMF or PDF without introducing additional numerical examples.

\section{Expectation of a Function of a Random Variable}

Let $Y = g(X)$. The expectation of $Y$ is:

\begin{itemize}
    \item \textbf{Discrete case}
    \[
    \mathbb{E}[g(X)] = \sum_x g(x)\, P(X = x).
    \]

    \item \textbf{Continuous case}
    \[
    \mathbb{E}[g(X)] = \int_{-\infty}^{+\infty} g(x)\, f_X(x)\, dx.
    \]
\end{itemize}

The expectation is computed directly with respect to the distribution of $X$.

\section{Linear Operation with Expectation}

\subsection{Linearity Property}

For random variables $X$ and $Y$, and constants $a$ and $b$:
\[
\mathbb{E}[aX + bY] = a\mathbb{E}[X] + b\mathbb{E}[Y].
\]

\subsubsection*{Justification}

\begin{enumerate}
    \item Apply the definition of expectation.
    \item Use linearity of summation or integration.
    \item Factor out constants.
\end{enumerate}

No independence assumption is required.

\section{Moments and Central Moments}

\subsection{nth Moment}

The \textbf{nth moment} of a random variable $X$ is
\[
\mathbb{E}[X^n].
\]

\subsection{Central Moments}

The \textbf{nth central moment} is
\[
\mathbb{E}[(X - \mu)^n], \quad \mu = \mathbb{E}[X].
\]

\subsection{Variance}

Variance is the second central moment:
\[
\mathrm{Var}(X) = \mathbb{E}[(X - \mu)^2].
\]

Using expansion,
\[
(X - \mu)^2 = X^2 - 2\mu X + \mu^2,
\]
which gives
\[
\mathrm{Var}(X) = \mathbb{E}[X^2] - (\mathbb{E}[X])^2.
\]

\subsection{Skewness}

Skewness corresponds to the third central moment and measures distribution asymmetry.

\subsection{Kurtosis}

Kurtosis corresponds to the fourth central moment and measures tail behavior.

\section*{End of Lecture 7 Scribe}

\end{document}
